% answers_short.tex —— 1977 年·山东A卷 · 简短答案速查
% 由 tools/make_short_answers.py 自 answers.tex 的 \gksolhead/\gkanswer 提取，
% 并与 1977/1977_数学真题_山东A_答案要点.md 逐题比对后生成，勿手改（改源 MD 后重跑）。

\begin{gknotice}
\textbf{使用说明}：本节只给最终答案便于快速核对，完整推导见解析版。答案为编写者独立推导，\textbf{非官方标准答案}。
\end{gknotice}

\gkpart{题目}
\begin{itemize}[leftmargin=2.6em]
  \item[\textbf{1.}] (1) $x\le3$ 时为 $3-x$，$x>3$ 时为 $x-3$；(2) $-\dfrac12$；(3) ① $x>0$；② $x\ge1$；(4) $y=-\dfrac12x+4$；(5) $F_1(-3,0)$，$F_2(3,0)$；(6) 正四棱台
  \item[\textbf{2.}] 见解析版
  \item[\textbf{3.}] 当 $0<\alpha<\dfrac{\pi}{4}$ 时，方程有两个不相等的实数根 $x=\dfrac{\cos\alpha\pm\sqrt{\cos2\alpha}}{2}$。
  \item[\textbf{4.}] (1) $r=\dfrac{2\pi-\alpha}{2\pi}R$，$h=\dfrac{\sqrt{4\pi\alpha-\alpha^2}}{2\pi}R$，$V=\dfrac{R^3(2\pi-\alpha)^2}{24\pi^2}\sqrt{4\pi\alpha-\alpha^2}$；(2) $V=1024\pi\approx3215\ \mathrm{cm^3}$
  \item[\textbf{5.}] $AB=\dfrac{\sqrt{15+6\sqrt3}}{3}\,a$。
  \item[\textbf{6.}] 每间长 $\dfrac{L}{4m}$、宽 $\dfrac{L}{2(m+1)}$ 时总面积最大，最大面积为 $\dfrac{L^2}{8(m+1)}$；此时 $m$ 间猪圈长的总和与宽的总和相等，均为 $\dfrac{L}{2}$。
\end{itemize}
